% What moves a claim's outcome in the model. It is the LaTeX drawing of the
% second mermaid diagram in docs/dispersion-model.md; build.py puts it in that
% diagram's place and stops when the two no longer match.
\begin{tikzpicture}[
  font=\small, >=Stealth,
  box/.style={draw, rounded corners=2pt, align=left, inner sep=5pt, minimum height=10mm, text width=#1},
  box/.default=74mm,
  flow/.style={->, semithick},
]
  \node[box] (u) {\textbf{paper} \ $u_p$\\how reproducible the paper is overall};
  \node[box, below=3mm of u] (w) {\textbf{claim} \ $w_{pc}$\\how hard this claim is};
  \node[box, below=3mm of w] (b) {\textbf{environment} \ $\beta_e$\\the shift a profile gives every paper};
  \node[box, below=3mm of b] (v) {\textbf{paper $\times$ environment} \ $v_{pe}$\\which papers a profile helps or hurts};
  \node[box, below=3mm of v] (eta) {\textbf{run} \ $\eta_{per}$\\a shock shared by all claims of one run};
  \node[box, dashed, below=9mm of eta] (s) {\textbf{spread of the run shock} \ $s_{pe}$\\set by the environment and by the paper:\\\mbox{$\log s_{pe}=\alpha_0+\alpha_e+\omega_p$}};

  \node[draw, rounded corners=8pt, fill=black!6, align=center, inner sep=6pt, text width=42mm, right=20mm of b] (y)
    {chance that claim $c$\\reproduces in run $r$\\[2pt]\mbox{$\operatorname{logit}\Pr(Y_{pcer}=1)$}};

  \foreach \source in {u, w, b, v, eta} \draw[flow] (\source.east) -- (y.west);
  \draw[flow] (s) -- (eta);
\end{tikzpicture}
